\begin{exercises}
  \recommended \item
    判断 \( \Re^3 \) 的下列每个子集是线性相关的还是线性无关的。
    \begin{exparts*}
      \partsitem \( \set{\colvec{1 \\ -3 \\ 5},
                    \colvec{2 \\ 2 \\ 4},
                    \colvec{4 \\ -4 \\ 14} }  \)
      \partsitem \( \set{\colvec{1 \\ 7 \\ 7},
                    \colvec{2 \\ 7 \\ 7},
                    \colvec{3 \\ 7 \\ 7} }  \)
      \partsitem \( \set{\colvec{0 \\ 0 \\ -1},
                    \colvec{1 \\ 0 \\ 4} }  \)
      \partsitem \( \set{\colvec{9 \\ 9 \\ 0},
                    \colvec{2 \\ 0 \\ 1},
                    \colvec{3 \\ 5 \\ -4},
                    \colvec{12 \\ 12 \\ -1} }  \)
    \end{exparts*}
    \begin{answer}
      对这里的每一小题，当子集是线性无关的时候你必须证明这一点，而当子集
      是线性相关的时候你必须给出一个相关关系的例子。
      \begin{exparts}
        \partsitem 它是线性相关的。
          考虑
          \begin{equation*}
             c_1\colvec{1 \\ -3 \\ 5}
             +c_2\colvec{2 \\ 2 \\ 4}
             +c_3\colvec{4 \\ -4 \\ 14}
             =\colvec{0 \\ 0 \\ 0}
          \end{equation*}
          就得到这个线性方程组。
          \begin{equation*}
            \begin{linsys}{3}
              c_1  &+  &2c_2  &+  &4c_3  &=  &0  \\
              -3c_1&+  &2c_2  &-  &4c_3  &=  &0  \\
              5c_1 &+  &4c_2  &+  &14c_3 &=  &0
            \end{linsys}
          \end{equation*}
          用高斯消元法
          \begin{equation*}
            \begin{amat}[r]{3}
              1  &2  &4  &0  \\
              -3 &2  &-4 &0  \\
              5  &4  &14 &0
            \end{amat}
            \grstep[-5\rho_1+\rho_3]{3\rho_1+\rho_2}
            \repeatedgrstep{(3/4)\rho_2+\rho_3}
            \begin{amat}[r]{3}
              1  &2  &4  &0  \\
              0  &8  &8  &0  \\
              0  &0  &0  &0
            \end{amat}
          \end{equation*}
          会得到一个自由变量，所以有无穷多解。
          作为相关关系的一个具体例子，我们可以取 $c_3$ 为，比如说，$1$。
          于是得到 \( c_2=-1 \) 和 \( c_1=-2 \)。
        \partsitem 它是线性相关的。
          这里出现的线性方程组
          \begin{equation*}
            \begin{amat}[r]{3}
              1  &2  &3  &0  \\
              7  &7  &7  &0  \\
              7  &7  &7  &0
            \end{amat}
            \grstep[-7\rho_1+\rho_3]{-7\rho_1+\rho_2}
            \repeatedgrstep{-\rho_2+\rho_3}
            \begin{amat}[r]{3}
              1  &2  &3   &0  \\
              0  &-7 &-14 &0  \\
              0  &0  &0   &0
            \end{amat}
          \end{equation*}
          有无穷多解。
          我们可以取 $c_3$ 为，比如说，$1$，再回代求出相应的 $c_2$ 和 $c_1$，
          从而得到一个特解。
        \partsitem 它是线性无关的。
          方程组
          \begin{equation*}
            \begin{amat}[r]{2}
              0  &1  &0  \\
              0  &0  &0  \\
              -1 &4  &0
            \end{amat}
            \grstep{\rho_1\leftrightarrow\rho_2}
            \repeatedgrstep{\rho_3\leftrightarrow\rho_1}
            \begin{amat}{2}
              -1 &4  &0  \\
              0  &1  &0  \\
              0  &0  &0
            \end{amat}
          \end{equation*}
          只有解 $c_1=0$ 与 $c_2=0$。
          （我们本来也可以凭观察得出答案\Dash 第二个向量显然不是第一个
          向量的倍数，反之亦然。）
        \partsitem 它是线性相关的。
          线性方程组
          \begin{equation*}
            \begin{amat}[r]{4}
              9  &2  &3  &12  &0  \\
              9  &0  &5  &12  &0  \\
              0  &1  &-4 &-1  &0
            \end{amat}
          \end{equation*}
          中未知量的个数多于方程的个数，因此用高斯消元法化简后最后必定
          至少有一个自由变量（不可能出现矛盾方程，因为该方程组是齐次的，
          从而至少有全零解）。
          为了展示出一个组合，我们可以作如下化简
          \begin{equation*}
            \grstep{-\rho_1+\rho_2}
            \grstep{(1/2)\rho_2+\rho_3}
            \begin{amat}[r]{4}
              9  &2  &3  &12  &0  \\
              0  &-2 &2  &0   &0  \\
              0  &0  &-3 &-1  &0
            \end{amat}
          \end{equation*}
          然后取，比如说，$c_4=1$。
          于是有 $c_3=-1/3$、$c_2=-1/3$ 以及 $c_1=-31/27$。
      \end{exparts}
    \end{answer}
  \recommended \item
    \( \polyspace_3 \) 的下列各子集中哪些是线性相关的，哪些是线性无关的？
    \begin{exparts}
      \partsitem \( \set{3-x+9x^2,5-6x+3x^2,1+1x-5x^2} \)
      \partsitem \( \set{-x^2,1+4x^2} \)
      \partsitem \( \set{2+x+7x^2,3-x+2x^2,4-3x^2} \)
      \partsitem \( \set{8+3x+3x^2,x+2x^2,2+2x+2x^2,8-2x+5x^2} \)
    \end{exparts}
    \begin{answer}
      在线性无关的情形，你必须证明它确实是线性无关的。否则，你必须展示出
      一个相关关系。（这里我们给出一个具体的相关关系，但也可能有别的。）
      \begin{exparts}
        \partsitem 这个集合是线性无关的。
          建立关系式
          \( c_1(3-x+9x^2)+c_2(5-6x+3x^2)+c_3(1+1x-5x^2)=0+0x+0x^2 \)
          就得到线性方程组
          \begin{equation*}
            \begin{amat}[r]{3}
              3  &5  &1  &0  \\
              -1 &-6 &1  &0  \\
              9  &3  &-5 &0
            \end{amat}
            \grstep[-3\rho_1+\rho_3]{(1/3)\rho_1+\rho_2}
            \repeatedgrstep{3\rho_2}
            \repeatedgrstep{-(12/13)\rho_2+\rho_3}
            \begin{amat}[r]{3}
              3  &5   &1        &0  \\
              0  &-13 &4        &0  \\
              0  &0   &-128/13  &0
            \end{amat}
          \end{equation*}
          只有唯一解：\( c_1=0 \)、\( c_2=0 \) 以及 \( c_3=0 \)。
        \partsitem 这个集合是线性无关的。
           我们可以直接从线性无关的定义出发，凭观察看出这一点。
           显然两者中任何一个都不是另一个的倍数。
        \partsitem 这个集合是线性无关的。
           线性方程组按如下方式化简
           \begin{equation*}
             \begin{amat}[r]{3}
               2  &3  &4  &0  \\
               1  &-1 &0  &0  \\
               7  &2  &-3 &0
             \end{amat}
             \grstep[-(7/2)\rho_1+\rho_3]{-(1/2)\rho_1+\rho_2}
             \repeatedgrstep{-(17/5)\rho_2+\rho_3}
             \begin{amat}[r]{3}
               2  &3    &4      &0  \\
               0  &-5/2 &-2     &0  \\
               0  &0    &-51/5  &0
             \end{amat}
           \end{equation*}
           由此看出只有解 $c_1=0$、$c_2=0$ 以及 $c_3=0$。
        \partsitem 这个集合是线性相关的。
           线性方程组
           \begin{equation*}
             \begin{amat}[r]{4}
               8  &0  &2  &8  &0  \\
               3  &1  &2  &-2 &0  \\
               3  &2  &2  &5  &0
             \end{amat}
           \end{equation*}
           在化简后必定最后至少有一个自由变量（变量的个数多于方程的
           个数，而且由于该方程组是齐次的，不可能出现矛盾方程）。
           我们可以把自由变量取作参数来描述解集。然后令参数取一个
           非零值，就得到一个非平凡的线性关系。
       \end{exparts}
      \end{answer}
  \item 判断下列每个集合在自然空间中是否线性无关。
    \begin{exparts*}
      \partsitem $\set{\colvec{1 \\ 2 \\ 0},
              \colvec{-1 \\ 1 \\ 0}}$
      \partsitem $\set{\rowvec{1 &3 &1}, \rowvec{-1 &4 &3}, \rowvec{-1 &11 &7}}$
      \partsitem $\set{\begin{mat}
                5 &4 \\
                1 &2
              \end{mat},
              \begin{mat}
                0 &0 \\
                0 &0
              \end{mat},
              \begin{mat}
                1 &0 \\
               -1 &4
              \end{mat}
           }$
    \end{exparts*}
    \begin{answer}
      \begin{exparts}
        \partsitem   自然向量空间是 $\Re^3$。
          建立方程
          \begin{equation*}
            \colvec{0 \\ 0 \\ 0}=c_1\colvec{1 \\ 2 \\ 0}
                                 +c_2\colvec{-1 \\ 1 \\ 0}
          \end{equation*}
          并考虑由此得到的齐次方程组。
          \begin{equation*}
             \begin{amat}{2}
               1  &-1 &0  \\
               2  &1  &0  \\
               0  &0  &0
             \end{amat}
             \grstep{-2\rho_1+\rho_2}
             \begin{amat}{2}
              1  &-1 &0  \\
              0  &3  &0  \\
              0  &0  &0
            \end{amat}
          \end{equation*}
          它有唯一解，即 $c_1=0$、$c_2=0$。
          所以它是线性无关的。
        \partsitem   自然向量空间是三个分量的行向量全体。
          方程
          \begin{equation*}
            \rowvec{0 &0 &0}=c_1\rowvec{1 &3 &1}
                            +c_2\rowvec{-1 &4 &3}
                            +c_3\rowvec{-1 &11 &7}
          \end{equation*}
          引出线性方程组
          \begin{align*}
            \begin{amat}{3}
              1  &-1  &-1  &0  \\
              3  &4   &11  &0  \\
              1  &3   &7   &0
           \end{amat}
           &\grstep[-\rho_1+\rho_3]{-3\rho_1+\rho_2}
           \begin{amat}{3}
             1  &-1  &-1  &0  \\
             0  &7   &14  &0  \\
             0  &4   &8   &0
          \end{amat}                                   \\
          &\grstep{(4/7)\rho_2+\rho_3}
          \begin{amat}{3}
            1  &-1  &-1  &0  \\
            0  &7   &14  &0  \\
            0  &0   &0   &0
         \end{amat}
       \end{align*}
       它有无穷多解，也就是说，不只是平凡解。
       \begin{equation*}
         \set{\colvec{c_1 \\ c_2 \\ c_3}
              =\colvec{-1 \\ -2 \\ 1}c_3
              \suchthat c_3\in\Re}
       \end{equation*}
       所以这个集合是线性相关的。
       一个相关关系来自取 $c_3=2$，此时得到 $c_1=-2$ 和 $c_2=-4$。
      \partsitem   不必建立方程组我们就能看出，集合中第二个元素是
        第一个元素的倍数（即第一个元素的 $0$ 倍）。
      \end{exparts}
    \end{answer}
  \recommended \item
    证明下列每个集合 \( \set{f,g} \) 在从 \( \Re^+ \) 到 \( \Re \) 的全体
    函数构成的向量空间中都是线性无关的。
    \begin{exparts}
      \partsitem \( f(x)=x \) 与 \( g(x)=1/x \)
      \partsitem \( f(x)=\cos(x) \) 与 \( g(x)=\sin(x) \)
      \partsitem \( f(x)=e^x \) 与 \( g(x)=\ln(x) \)
    \end{exparts}
    \begin{answer}
      设 $\map{Z}{\Re^+}{\Re}$ 为零函数 $Z(x)=0$，它是所讨论的向量空间中
      的加法单位元。
      \begin{exparts}
        \partsitem 这个集合是线性无关的。
          考虑 \( c_1\cdot f(x)+c_2\cdot g(x)=Z(x) \)。
          代入 \( x=1 \) 和 \( x=2 \) 就得到线性方程组
          \begin{equation*}
            \begin{linsys}{2}
              c_1\cdot 1  &+  &c_2\cdot 1     &=  &0  \\
              c_1\cdot 2  &+  &c_2\cdot (1/2) &=  &0
            \end{linsys}
          \end{equation*}
          它有唯一解 \( c_1=0 \)、\( c_2=0 \)。
        \partsitem 这个集合是线性无关的。
          考虑 \( c_1\cdot f(x)+c_2\cdot g(x)=Z(x) \)，
          代入 \( x=\pi \) 和 \( x=\pi/2 \) 得到
          \begin{equation*}
            \begin{linsys}{2}
              c_1\cdot (-1)  &+  &c_2\cdot 0     &=  &0  \\
              c_1\cdot 0     &+  &c_2\cdot 1     &=  &0
            \end{linsys}
          \end{equation*}
          它显然给出 \( c_1=0 \)、\( c_2=0 \)。
        \partsitem 这个集合也是线性无关的。
          考虑 \( c_1\cdot f(x)+c_2\cdot g(x)=Z(x) \)
          并代入 \( x=1 \) 和 \( x=e \)
          \begin{equation*}
            \begin{linsys}{2}
              c_1\cdot e    &+  &c_2\cdot 0     &=  &0  \\
              c_1\cdot e^e  &+  &c_2\cdot 1     &=  &0
            \end{linsys}
          \end{equation*}
          得到 \( c_1=0 \) 与 \( c_2=0 \)。
      \end{exparts}
     \end{answer}
  \recommended \item
    实变量实值函数空间中，下列各子集哪些是线性相关的，哪些是线性
    无关的？
    （我们对一些常值函数作了缩写；~例如，第一小题中的 `$2$' 就代表
     常值函数 $f(x)=2$。）
    \begin{exparts*}
      \partsitem \( \set{2,4\sin^2(x),\cos^2(x)} \)
      \partsitem \( \set{1,\sin(x),\sin(2x)} \)
      \partsitem \( \set{x,\cos(x)} \)
      \partsitem \( \set{(1+x)^2,x^2+2x,3} \)
      \partsitem \( \set{\cos(2x),\sin^2(x),\cos^2(x)} \)
      \partsitem \( \set{0,x,x^2} \)
    \end{exparts*}
    \begin{answer}
      在每一种情形，如果集合是线性无关的，那么你必须证明这一点；如果它是
      线性相关的，那么你必须展示出一个相关关系。
      \begin{exparts}
        \partsitem 这个集合是线性相关的。
          熟知的关系式 $\sin^2(x)+\cos^2(x)=1$ 表明，
          $2=c_1\cdot(4\sin^2(x))+c_2\cdot(\cos^2(x))$ 在
          $c_1=1/2$ 和 $c_2=2$ 时成立。
        \partsitem 这个集合是线性无关的。
          考虑关系式
          $c_1\cdot 1+c_2\cdot\sin(x)+c_3\cdot\sin(2x)=0$
          （其中 `$0$' 指零函数）。
          取三个适当的点，比如 $x=\pi$、$x=\pi/2$、$x=\pi/4$，
          就得到方程组
          \begin{equation*}
            \begin{linsys}{3}
               c_1  &   &                &   &      &=  &0  \\
               c_1  &+  &c_2             &   &      &=  &0  \\
               c_1  &+  &(\sqrt{2}/2)c_2 &+  &c_3   &=  &0
            \end{linsys}
          \end{equation*}
          其唯一解是
          $c_1=0$、$c_2=0$ 以及 $c_3=0$。
        \partsitem 凭观察可知，这个集合是线性无关的。
          任何形如 $\cos(x)=c\cdot x$ 的相关关系都不可能成立，因为余弦
          函数不是恒等函数的倍数
          （我们在这里应用 \nearbycorollary{cor:LDMeansLC}）。
        \partsitem 凭观察我们发现存在一个相关关系。
          因为 $(1+x)^2=x^2+2x+1$，所以
          $c_1\cdot(1+x)^2+c_2\cdot(x^2+2x)=3$ 在
          $c_1=3$ 和 $c_2=-3$ 时成立。
        \partsitem 这个集合是线性相关的。
          看出这一点最简单的方法是想起三角关系式
          $\cos^2(x)-\sin^2(x)=\cos(2x)$。
          （\textit{注}：
          想不起来这个关系式的人，若试着代入一些 $x$ 值，只会始终得不到
          通向唯一解的方程组，也就永远无法得出该集合线性无关的结论。
          当然，这个人可能会怀疑是不是恰好没试对那组 $x$ 值，但试过
          几次之后，大多数人就会转而去寻找相关关系。）
        \partsitem 这个集合是线性相关的，因为它包含该向量空间中的
          零对象，即零多项式。
      \end{exparts}
     \end{answer}
  \item
    方程 \( \sin^2(x)/\cos^2(x)=\tan^2(x) \) 是否表明，函数集合
    \( \set{\sin^2(x),\cos^2(x),\tan^2(x)} \) 是全体实值函数之集合的
    一个线性相关子集？这里全体实值函数的定义域是由介于
    \( -\pi/2 \) 与 \( \pi/2) \) 之间的实数构成的区间
    \( (-\pi/2..\pi/2) \)。
    \begin{answer}
      不是，那个方程并不是一个线性关系。
      事实上这个集合是线性无关的，取 \( x \) 为 \( 0 \)、\( \pi/6 \)
      和 \( \pi/4 \) 所得到的方程组即可说明这一点。
    \end{answer}
  \item
    向量空间 $\Re^3$ 的子集 $xy$ 平面是线性无关的吗？
    \begin{answer}
      不是。
      该平面中有两个元素，第二个是第一个的倍数。
      \begin{equation*}
        \colvec{1  \\ 0 \\ 0},\,\colvec{2 \\ 0 \\ 0}
      \end{equation*}
      （答案为“不是”的另一个理由是，零向量是该平面的元素，而包含零
      向量的集合都不可能是线性无关的。）
    \end{answer}
  \recommended \item
    证明阶梯形矩阵的非零行构成一个线性无关的集合。
    \begin{answer}
      我们已经证明过这一点：线性组合引理及其推论指出，在阶梯形矩阵中，
      任何非零行都不是其余各行的线性组合。
    \end{answer}
  \item
     \begin{exparts}
       \partsitem 证明：如果集合 \( \set{\vec{u},\vec{v},\vec{w}} \)
          是线性无关的，那么集合
          \( \set{\vec{u},\vec{u}+\vec{v},\vec{u}+\vec{v}+\vec{w}} \)
          也是线性无关的。
       \partsitem  \( \set{\vec{u},\vec{v},\vec{w}} \) 的线性无关性或
         线性相关性与
         \( \set{\vec{u}-\vec{v},\vec{v}-\vec{w},\vec{w}-\vec{u}} \) 的
         无关性或相关性之间有什么关系？
     \end{exparts}
     \begin{answer}
       \begin{exparts}
         \partsitem 假设
           \( \set{\vec{u},\vec{v},\vec{w}} \) 是
           线性无关的，于是任何关系式
           $d_0\vec{u}+d_1\vec{v}+d_2\vec{w}=\zero$ 都导致
           $d_0=0$、$d_1=0$ 以及 $d_2=0$ 的结论。

           考虑关系式
           \( c_1(\vec{u})+c_2(\vec{u}+\vec{v})+c_3(\vec{u}+\vec{v}+\vec{w})
           =\zero \)。
           将其改写，得到
           \( (c_1+c_2+c_3)\vec{u}+(c_2+c_3)\vec{v}+(c_3)\vec{w}=\zero \)。
           取 $d_0$ 为 $c_1+c_2+c_3$，取 $d_1$ 为 $c_2+c_3$，
           取 $d_2$ 为 $c_3$，就得到这个方程组。
           \begin{equation*}
             \begin{linsys}{3}
               c_1  &+  &c_2  &+  &c_3  &=  &0  \\
                    &   &c_2  &+  &c_3  &=  &0  \\
                    &   &     &   &c_3  &=  &0
              \end{linsys}
            \end{equation*}
            结论：~各 $c$ 全为零，所以该集合是线性
            无关的。
         \partsitem 第二个集合是线性相关的
            \begin{equation*}
              1\cdot(\vec{u}-\vec{v})
              +1\cdot(\vec{v}-\vec{w})
              +1\cdot(\vec{w}-\vec{u})
              =\zero
            \end{equation*}
            无论第一个集合是否线性无关。
       \end{exparts}
     \end{answer}
  \item
    \nearbyexample{ex:EmSetLI} 表明空集是线性无关的。
    \begin{exparts}
      \partsitem 单元素集何时是线性无关的？
      \partsitem 两个元素的集合又如何？
    \end{exparts}
    \begin{answer}
      \begin{exparts}
        \partsitem 单元素集 $\set{\vec{v}}$ 是线性无关的
          当且仅当 $\vec{v}\neq\zero$。
          对于“若”的方向，当 $\vec{v}\neq\zero$ 时，考虑
          关系式
          \( c\cdot\vec{v}=\zero \) 并注意唯一解
          是平凡的：~$c=0$，就可以应用
          \nearbylemma{le:LDIffANonTrivLinRel}。
          对于“仅当”的方向，只需回忆
          \nearbyexample{ex:SetWithZeroVecLD}
          表明 $\set{\zero}$ 是线性相关的，因此如果集合
          $\set{\vec{v}}$ 是线性无关的，那么 $\vec{v}\neq\zero$。

          （\textit{注}：
          另一种回答是说，这是
          \nearbylemma{lm:AddVecLiSetIsLiIffVecNotInSpan}
          在 \( S=\emptyset \) 时的特殊情形。）
        \partsitem 两个元素的集合是线性无关的
          当且仅当其中任何一个元素都不是另一个的倍数
          （注意，如果其中一个是零向量，那么它是另一个的倍数）。
          等价的说法：~一个集合是线性相关的当且仅当
          其中一个元素是另一个的倍数。

          证明很容易。
          集合 $\set{\vec{v}_1,\vec{v}_2}$ 是线性相关的当且仅当
          存在关系式 $c_1\vec{v}_1+c_2\vec{v}_2=\zero$，
          其中 $c_1\neq 0$ 或 $c_2\neq 0$（或两者都成立）。
          这当且仅当 $\vec{v}_1=(-c_2/c_1)\vec{v}_2$
          或 $\vec{v}_2=(-c_1/c_2)\vec{v}_1$（或两者都成立）时成立。
       \end{exparts}
     \end{answer}
  \item
    在任何向量空间 \( V \) 中，空集都是线性无关的。
    那么 \( V \) 本身呢？
    \begin{answer}
      这个集合是线性相关的，因为它包含零向量。
    \end{answer}
  \item
    证明：如果 \( \set{\vec{x},\vec{y},\vec{z}} \) 是线性
    无关的，那么它的所有真子集：~\( \set{\vec{x},\vec{y}} \)、
    \( \set{\vec{x},\vec{z}} \)、\( \set{\vec{y},\vec{z}} \)、
    \( \set{\vec{x}} \)、\( \set{\vec{y}} \)、\( \set{\vec{z}} \)
    以及 \( \set{} \) 也都是线性无关的。
    反过来（“仅当”）也成立吗？
    \begin{answer}
      \nearbylemma{le:SubsetPreserveLI} 给出了“若”的那一半。
      逆命题（“仅当”的陈述）不成立。
      举个例子：考虑向量空间 \( \Re^2 \) 和
      这些向量。
      \begin{equation*}
         \vec{x}=\colvec{1 \\ 0},\quad
         \vec{y}=\colvec{0 \\ 1},\quad
         \vec{z}=\colvec{1 \\ 1}
      \end{equation*}
    \end{answer}
  \item
    \begin{exparts}
      \partsitem 证明：这个
        \begin{equation*}
          S=\set{\colvec{1 \\ 1 \\ 0},\colvec{-1 \\ 2 \\ 0}}
        \end{equation*}
        是 \( \Re^3 \) 的一个线性无关子集。
      \partsitem 通过求出给出线性关系的 \( c_1 \) 和 \( c_2 \)，
        证明
        \begin{equation*}
          \colvec{3 \\ 2 \\ 0}
        \end{equation*}
        在 $S$ 的张成中。
        \begin{equation*}
          c_1\colvec{1 \\ 1 \\ 0}
          +c_2\colvec{-1 \\ 2 \\ 0}
          =\colvec{3 \\ 2 \\ 0}
        \end{equation*}
        证明数对 \( c_1,c_2 \) 是唯一的。
      \partsitem 假设 \( S \) 是向量空间的一个子集，且
        \( \vec{v} \) 属于 \( \spanof{S} \)，即 \( \vec{v} \) 是
        $S$ 中向量的一个线性组合。
        证明：如果 \( S \) 是线性无关的，那么 $S$ 中向量相加等于
        \( \vec{v} \) 的线性组合是唯一的（也就是说，唯一性是在重新
        排列次序以及添加或删去形如 \( 0\cdot\vec{s} \) 的项的意义下
        而言的）。
        因此，作为张成集的 \( S \) 在这个强意义下是极小的：
        \( \spanof{S} \) 中的每个向量都是 $S$ 的元素以最少
        次数\Dash 仅一次\Dash 构成的组合。
      \partsitem
        证明：当 \( S \) 不是线性无关的时候，可能发生不同的线性组合
        求和得到同一个向量的情形。
    \end{exparts}
    \begin{answer}
      \begin{exparts}
        \partsitem 由
          \begin{equation*}
            c_1\colvec{1 \\ 1 \\ 0}
            +c_2\colvec{-1 \\ 2 \\ 0}
            =\colvec{0 \\ 0 \\ 0}
          \end{equation*}
          得到的线性方程组有唯一解 \( c_1=0 \) 和 \( c_2=0 \)。
        \partsitem 由
          \begin{equation*}
            c_1\colvec{1 \\ 1 \\ 0}
            +c_2\colvec{-1 \\ 2 \\ 0}
            =\colvec{3 \\ 2 \\ 0}
          \end{equation*}
          得到的线性方程组有唯一解 \( c_1=8/3 \) 和 \( c_2=-1/3 \)。
        \partsitem 假设 \( S \) 是线性无关的。
          再假设我们同时有 $\vec{v}=c_1\vec{s}_1+\dots+c_n\vec{s}_n$
          和 $\vec{v}=d_1\vec{t}_1+\dots+d_m\vec{t}_m$
          （其中这些向量都是 $S$ 的元素）。
          于是，
          \begin{equation*}
            c_1\vec{s}_1+\dots+c_n\vec{s}_n
            =\vec{v}
            =d_1\vec{t}_1+\dots+d_m\vec{t}_m
          \end{equation*}
          可以如下改写。
          \begin{equation*}
            c_1\vec{s}_1+\dots+c_n\vec{s}_n
            -d_1\vec{t}_1-\dots-d_m\vec{t}_m
            =\zero
          \end{equation*}
          可能有些 $\vec{s}$ 与某些 $\vec{t}$ 相等；
          我们可以把相应的系数合并
          （即，如果 $\vec{s}_i=\vec{t}_j$，那么
          $\cdots+c_i\vec{s}_i+\dots-d_j\vec{t}_j-\cdots$ 可以改写
          为 $\cdots+(c_i-d_j)\vec{s}_i+\cdots$）。
          这个等式是集合 $S$ 的（在合并完成之后的）不同元素之间的
          一个线性关系。
          我们已假设 $S$ 是线性无关的，所以所有系数都为零。
          如果 $i$ 使得 $\vec{s}_i$ 不等于任何 $\vec{t}_j$，
          那么 $c_i$ 为零。
          如果 $j$ 使得 $\vec{t}_j$ 不等于任何 $\vec{s}_i$，
          那么 $d_j$ 为零。
          在最后一种情形，我们有 $c_i-d_j=0$，从而 $c_i=d_j$。

          因此，原来的两个和是相同的，也许除去一些
          $0\cdot\vec{s}_i$ 或 $0\cdot\vec{t}_j$ 项，这些项可以忽略。
        \partsitem
          这个集合不是线性无关的：
          \begin{equation*}
            S=\set{\colvec{1 \\ 0},\colvec{2 \\ 0}}\subset\Re^2
          \end{equation*}
          而这两个线性组合给出同样的结果
          \begin{equation*}
            \colvec{0 \\ 0}=2\cdot\colvec{1 \\ 0}-1\cdot\colvec{2 \\ 0}
                            =4\cdot\colvec{1 \\ 0}-2\cdot\colvec{2 \\ 0}
          \end{equation*}
          因此，一个线性相关的集合的和可能不是互不相同的。

          事实上，一个更强的论断成立：~如果一个集合是线性相关的，那么它
          必定具有这样的性质：存在两个不同的线性组合求和得到同一个向量。
          简言之，若有 \( c_1\vec{s}_1+\dots+c_n\vec{s}_n=\zero \)，
          则将关系式两边乘以二就得到另一个关系式。
          如果第一个关系式是非平凡的，那么第二个也是。
      \end{exparts}
    \end{answer}
  \item  \label{exer:PolyZeroFcnOnlyIfZeroPol}
     证明：一个多项式导出零函数当且仅当
     它是零多项式。
     （\textit{评注}：
     这个问题不是线性代数的问题，但我们经常用到这个结论。
     多项式以自然的方式导出一个
     函数：~$x\mapsto c_nx^n+\dots+c_1x+c_0$。）
     \begin{answer}
       在这个“当且仅当”的陈述中，“若”的那一半是清楚的\Dash 如果
       多项式是零多项式，那么由该多项式的作用导出的函数必定是零
       函数 $x\mapsto 0$。
       对于“仅当”，我们写 $p(x)=c_nx^n+\dots+c_0$。
       代入零，$p(0)=0$，得到 $c_0=0$。
       取导数再代入零，$p^\prime(0)=0$，得到
       $c_1=0$。
       类似地，我们得到每个 $c_i$ 都为零，于是 $p$ 是零
       多项式。
     \end{answer}
  \item
    回到 1.2 节，重新定义点、直线、平面
    以及其他线性曲面，以避免退化的情形。
    \begin{answer}
      本节的工作提示我们，应当把 \( n \) 维
      非退化线性曲面定义为 \( n \) 个向量组成的线性无关集合的
      张成。
    \end{answer}
  \item
    \begin{exparts}
      \partsitem 证明 \( \Re^2 \) 中任意四个向量构成的集合都是
         线性相关的。
      \partsitem 对任意五个向量的集合这也成立吗？
         任意三个呢？
      \partsitem $\Re^2$ 的线性无关子集最多能有多少个元素？
    \end{exparts}
    \begin{answer}
      \begin{exparts}
        \partsitem 对任意的 $a_{1,1}$、\ldots、$a_{2,4}$，
          \begin{equation*}
            c_1\colvec{a_{1,1} \\ a_{2,1}}
            +c_2\colvec{a_{1,2} \\ a_{2,2}}
            +c_3\colvec{a_{1,3} \\ a_{2,3}}
            +c_4\colvec{a_{1,4} \\ a_{2,4}}
            =\colvec{0 \\ 0}
          \end{equation*}
          给出线性方程组
          \begin{equation*}
             \begin{linsys}{4}
              a_{1,1}c_1 &+ &a_{1,2}c_2 &+ &a_{1,3}c_3 &+ &a_{1,4}c_4 &= &0  \\
              a_{2,1}c_1 &+ &a_{2,2}c_2 &+ &a_{2,3}c_3 &+ &a_{2,4}c_4 &= &0
             \end{linsys}
          \end{equation*}
        它有无穷多解（高斯消元法至少留下
        两个自由变量）。
        因此，$\Re^2$ 中给定的这些元素之间存在非平凡的线性关系。
      \partsitem 任意五个向量的集合是某个四向量集合的超集，
        所以是线性相关的。

